\documentclass[11pt,a4paper]{exam}

\usepackage[utf8]{inputenc}
\usepackage[T1]{fontenc}
\usepackage[english,main=french]{babel}
\usepackage{amsmath, amssymb}
\usepackage{geometry}
\usepackage{listings}
\usepackage{xcolor}

\geometry{margin=2cm}

\lstset{
  basicstyle=\ttfamily\small,
  keywordstyle=\bfseries,
  columns=fullflexible,
  keepspaces=true,
  frame=single,
  showstringspaces=false,
  tabsize=2
}

\firstpageheader{BTS CIEL1}{Python}{Date : 02/12/2025}
\runningheader{Python}{}{Page \thepage/\numpages}

\begin{document}

\begin{center}
  \Large\textbf{Algorithmes et structures de données en Python}\\[0.5em]
  \normalsize Durée : 1h \\
  \vspace{1em}
  \textbf{Nom :} \rule{6cm}{0.4pt} \hfill
  \textbf{Prénom :} \rule{6cm}{0.4pt}\\[1em]
\end{center}

\vspace{1em}

\section*{Exercice 1 — QCM (10 questions)}

Cochez la bonne réponse pour chaque question. \textbf{Une seule} réponse correcte par question.

\begin{questions}

\question[1] Python est un langage :
\begin{checkboxes}
  \choice Seulement dédié au web
  \choice Interprété et polyvalent (GPL)
  \choice Réservé aux data scientists
  \choice Basé sur une compilation stricte
\end{checkboxes}

\question[1] Quelle structure permet de stocker une séquence modifiable d'éléments ?
\begin{checkboxes}
  \choice Tuple
  \choice Dictionnaire
  \choice Liste
  \choice String
\end{checkboxes}

\question[1] Que signifie « immuable » pour un tuple ?
\begin{checkboxes}
  \choice Les éléments ne peuvent pas être modifiés après la création
  \choice On ne peut jamais accéder aux éléments
  \choice Les éléments doivent être de même type
  \choice Il ne peut pas contenir plus de 3 éléments
\end{checkboxes}

\question[1] Dans un dictionnaire Python, l’accès à une valeur se fait via :
\begin{checkboxes}
  \choice Un indice entier
  \choice Une clé unique
  \choice Une position aléatoire
  \choice Une étiquette numérique
\end{checkboxes}

\question[1] L’instruction \texttt{in} dans un dictionnaire vérifie :
\begin{checkboxes}
  \choice L’existence d’une valeur
  \choice L’existence d’une clé
  \choice Le type du dictionnaire
  \choice Le nombre d’éléments
\end{checkboxes}

\question[1] Quelle structure ne permet \textbf{pas} les doublons ?
\begin{checkboxes}
  \choice Liste
  \choice String
  \choice Ensemble (set)
  \choice Tuple
\end{checkboxes}

\question[1] En Python, la fonction qui permet d’obtenir une entrée utilisateur est :
\begin{checkboxes}
  \choice \texttt{scanf()}
  \choice \texttt{read()}
  \choice \texttt{input()}
  \choice \texttt{enter()}
\end{checkboxes}

\newpage

\question[1] La borne supérieure de \texttt{range(a, b)} est :
\begin{checkboxes}
  \choice Incluse
  \choice Toujours 0
  \choice Exclue
  \choice Dépend de l’indentation
\end{checkboxes}

\question[1] On veut tester la présence d’un élément rapidement. Quel type est le plus adapté ?
\begin{checkboxes}
  \choice Liste
  \choice Tuple
  \choice Set
  \choice String
\end{checkboxes}

\question[1] Quel type de structure est le plus adapté pour stocker et modifier les informations d’un étudiant (nom, prénom, âge, formation) ?
\begin{checkboxes}
  \choice Tuple
  \choice Liste
  \choice Dictionnaire
  \choice Ensemble
\end{checkboxes}

\end{questions}

\section*{Exercice 2 — Programme Python à trous}

\begin {questions}

\question[2] Complétez le programme Python (Micro:bit) ci-dessous :

\begin{lstlisting}[language=Python]
from microbit import *

import ________

MAX = ________

number_to_guess = random.________(1, MAX)
answer = 0

display.________()

while True:
    display.show(question_mark_image)
            
    if ________.is_pressed():
        answer = (answer + 1) % (MAX + 1)
        display.scroll(answer, delay=50, wait=True)

    if ________.was_pressed():
        if answer > number_to_guess:
            display.show(Image.SAD)
            display.show(Image.LESS)
            answer = 0
        if answer < number_to_guess:
            display.show(Image.SAD)
            display.show(Image.PLUS)
        else:
            display.show(Image.HAPPY)
            answer = 0
        sleep(_______)

\end{lstlisting}

\question[0.5] Expliquez l'instruction suivante : \verb|(answer + 1) % (MAX + 1)|

\end {questions}

\newpage

\section*{Exercice 3 — Traduction de C vers Python}

\begin{questions}
\question On donne le programme C suivant :

\begin{lstlisting}[language=C]
#include <stdio.h>

int main(int argc, char const *argv[])
{
    int e = 0;
    int nbOcc = 0;
    int posOcc = -1;

    int n = -1;
    int tab[100] = {0};

    for (int i = 0; i < n; i++) scanf("%d", &tab[i]);

    printf("Quel element souhaitez-vous rechercher :");
    scanf("%d", &e);

    for (int i = 0; i < n; i++) {
        if (tab[i] == e) {
            nbOcc++;
            posOcc = i;
        }
    }

    printf("Nombre d'occurences : %d, position : %d \n", nbOcc, posOcc);
}
\end{lstlisting}

\begin{parts}
  \part[0.5] Expliquez en quelques mots ce que fait ce programme.
  \vspace{3cm}

  \part[2] Transcrivez ce programme en langage \textbf{Python} ci-dessous :

  \vspace{0.5cm}
  \begin{solutionbox}{8cm}
  \end{solutionbox}
\end{parts}
\end{questions}

\newpage

\section*{Exercice 4 — Programme Micro:bit}

La carte \textbf{micro:bit} possède une \textbf{matrice de 25 LED} (5 colonnes × 5 lignes).  

Chaque LED est identifiée par ses coordonnées :

\begin{itemize}
    \item \texttt{x} : numéro de la colonne (0 à 4)
    \item \texttt{y} : numéro de la ligne (0 à 4)
\end{itemize}

Pour lire ou modifier l'état d'une LED, on peut utiliser :

\begin{itemize}
    \item \verb|display.get_pixel(x, y)| : renvoie l'état actuel (0 = éteinte, 9 = allumée)
    \item \verb|display.set_pixel(x, y, valeur)| : modifie l'état
\end{itemize}

\begin{questions}

\question[5] \textbf{Ecrivez un programme Python qui :}
\begin{itemize}
    \item parcourt les \textbf{25 LED} de la matrice (de gauche à droite et de haut en bas),
    \item inverse l'état de chaque LED :
    \begin{itemize}
        \item si elle est éteinte → elle s'allume,
        \item si elle est allumée → elle s'éteint,
    \end{itemize}
    \item passe ensuite à la LED suivante,
    \item recommence indéfiniment pour former une animation de serpentin.
\end{itemize}

\vspace{0.5cm}
\noindent\textbf{Votre programme Python :}

\begin{solutionbox}{15cm}
\end{solutionbox}

\end{questions}


\end{document}
